\begin{abstract}
% AI agents perform well on individual software engineering tasks but face architectural bottlenecks on long-horizon tasks requiring many interdependent changes across files, while existing multi-agent approaches primarily split roles without achieving true parallel execution. 
% However, the core challenges for multi-agent collaboration on such tasks, including workspace isolation, dependency scheduling, output integration, and correctness verification, which already have mature solutions in software engineering practice: git worktrees, import graphs, git merge, and test suites. 
AI agents have become increasingly capable at isolated software engineering (SWE) tasks such as resolving issues on Github.
Yet long-horizon tasks involving multiple interdependent subtasks still pose challenges both with respect to accuracy, and with respect to timely completion.
A natural approach to solving these long-horizon tasks in a timely manner is asynchronous multi-agent collaboration, where multiple agents work on different parts of the task at the same time.
But effective application of multi-agent systems has proven surprisingly difficult: concurrent edits by multiple agents interfere with each other, dependencies are difficult to synchronize, and combining partial progress into a coherent whole is challenging.
% In the shared-artifact and long-horizon regime, software engineering is a natural testbed, and 
On the other hand, human developers have long relied on mature collaboration infrastructure to manage these challenges in large software projects.
Inspired by these collaboration primitives, we introduce \underline{\textit{C}}entralized \underline{\textit{A}}synchronous \underline{\textit{I}}solated \underline{\textit{D}}elegation (\ourmethod), a structured multi-agent coordination paradigm grounded in three core SWE primitives: centralized task delegation, asynchronous execution, and isolated workspaces. \ourmethod constructs dependency-aware task plans through a central manager, executes subtasks concurrently in isolated workspaces, and consolidates progress via structured integration with executable test-based verification. In empirical evaluation, we find that \ourmethod improves accuracy over single-agent baselines by 26.7\% absolute on paper reproduction tasks (PaperBench) and 14.3\% on Python library development tasks (Commit0).
Through systematic analysis, we find that branch-and-merge is a central coordination mechanism for multi-agent collaboration, and that SWE primitives such as \texttt{git worktree}, \texttt{git commit}, and \texttt{git merge} enable it to be realized in a reliable and executable manner.
% \begin{figure*}[t]
%     \centering
%     \includegraphics[width=\textwidth]{pdfs/multi-agent-teaser.pdf}
%     \caption{
%        \textbf{Overview of \ourmethod Workflow.} The Manager analyzes the codebase (commit0) or the main contribution of the paper (paperbench), builds a dependency graph, decomposes tasks into parallelizable groups, and creates isolated git worktrees. In the asynchronous loop, engineers independently implement, self-verify, and commit, upon any engineer's completion, the Manager merges to main and dynamically updates the task delegation plan before reassigning the next task. After the asynchronous loop, the manager does a final review before submitting to evaluate.
%        % \jy{update} \gncomment{Some comments: (1) I didn't really understand ``B'' here, maybe just ``A'' is sufficient. (2) I don't think you need the OpenHands logo, it's confusing have two different logos to represent agents, so probably the panda robot is fine. (3) the logos in ``Manager analyzes tasks'' were not very clear. Maybe they were meant to be different task types, but I didn't get that from the figure. Maybe you could just have one logo to represent ``task''. (4) Make sure that you're consistent in title case. There are some typos like ``Isolated ask delegation''.}
%     }
%     \label{fig:teaser}
% \end{figure*}

\begin{figure}[h]
    \centering
    \includegraphics[width=0.95\textwidth]{pdfs/multi-agent-teaser.pdf}
    \caption{
       \textbf{Overview of \ourmethod Workflow.} The Manager explores the SWE tasks, builds a dependency graph to decompose tasks into parallelizable groups, and creates isolated git worktrees for every onboarded engineer. In the asynchronous loop, engineers independently implement, self-verify, and make a commit. Upon any engineer's completion, the Manager merges to main and dynamically updates the task delegation plan before reassigning the next task. After the asynchronous loop, the manager does a final review before submitting the final product. }
    \label{fig:teaser}
\end{figure}
       % \jy{manager does the git merge and assign the "resolve conflict task" to the engineer}
       % \jy{update} \gncomment{Some comments: (1) I didn't really understand ``B'' here, maybe just ``A'' is sufficient. (2) I don't think you need the OpenHands logo, it's confusing have two different logos to represent agents, so probably the panda robot is fine. (3) the logos in ``Manager analyzes tasks'' were not very clear. Maybe they were meant to be different task types, but I didn't get that from the figure. Maybe you could just have one logo to represent ``task''. (4) Make sure that you're consistent in title case. There are some typos like ``Isolated ask delegation''.}



% Through systematic analysis, we further find that worktree isolation is essential for execution stability and the degree of parallelism must match task modularity and manager's task delegation ability.
\end{abstract}